\section{Evola on politics, morals, philosophy}

These are the first two questions and answers from the interview published in ``La Nation Europeenne''.

\textbf{Q}. \emph{Do you believe that there exists a connection between philosophy and politics? Can a philosophy influence an undertaking of national or European political reconstruction?}

\textbf{A}. I don't believe that philosophy in the strictly theoretical sense of the term can have any influence over politics. In order to exercise an influence, it needs to embody an ideology or a conception of the world. It is what happened, for example, with the Enlightenment, with Marxist dialectical materialism, and with certain philosophic conceptions that were incorporated in the conception of the world of German National Socialism. In general, the era of grand philosophical systems is over, now there only exist illegitimate and mediocre philosophies. In one of my early works from my philosophical period, I had included these words of Jules Lachelier actually, Jules Lagneau : ``Modern philosophy is a reflexion that has ended by recognising its own impotence and the necessity of action that arises from within.'' The proper domain of an action of this type has a meta-philosophical character. Hence, the transition that can be observed in my books which speak not of ``philosophy'', but of ``metaphysics'', of worldviews, and of traditional doctrines.

\textbf{Q}. \emph{Do you think that morals and ethics are synonymous and should share a philosophy foundation?}

\textbf{A}. It is possible to establish a distinction, if by ``morals'' you mean customs and by ``ethics'', a philosophical discipline (which is called ``moral philosophy''). It seems to me, it is illusory to demand an absolute philosophical foundation for any ethics or morals. Without reference to something transcendent, morals can only have a relative, contingent, and social importance and can not overcome the criticisms of individualism, existentialism, and nihilism. I demonstrated it in my book ``Ride the Tiger'' in the chapter ``In the World where God is Dead''. In this chapter I also dealt with the problematic posed by Nietzsche and existentialism.

\flrightit{Posted on 2007-12-31 by Cologero}

